\exercisesection{附录~II}

¶ 1) 设 X 是局部连通拓扑空间，\(\mathcal{C}(\mathrm{X})\) 是 X 上实值连续函数组成的空间，\(d\) 是整数 \(\geq 0\) 。设 \(\mathrm{F} \in \mathcal{C}(\mathrm{X})[\mathrm{T}]\) 是次数为 \(d\) 、系数属于 \(\mathcal{C}(\mathrm{X})\) 的首一多项式：

\[
\mathrm{F} = \mathrm{T} ^ {d} + \mathrm{T} ^ {d - 1} f _ {1} + \dots + f _ {d}, \quad f _ {i} \in \mathcal {C} (\mathrm{X}).
\]

通过下式把 F 等同于 \(R \times X\) 上的一个函数：

\[
\mathrm{F} (t, x) = t ^ {d} + t ^ {d - 1} f _ {1} (x) + \dots + f _ {d} (x) \text {if} t \in \mathbf {R}, x \in \mathrm{X}.
\]

设 \(\Delta \in \mathcal{C}(\mathrm{X})\) 是多项式 F 的判别式（《代数学》第四章 §1 第 10 节）。

若 U 是 X 的开子集，用 \(Z_{U}\) 表示由 \((t,x)\) （其中 \(t \in R\) 且 \(x \in U\) ）组成且满足 \(F(t,x) = 0\) 的集合；这是 \(R \times U\) 的闭子集。

\begin{enumerate}
\def\labelenumi{\alph{enumi})}
\item
  证明投影 \(\mathrm{pr}_2: \mathrm{Z}_{\mathrm{U}} \to \mathrm{U}\) 是真的（《一般拓扑学》第一章 §10）。
\item
  假设 U 连通，且 \(\Delta(x) \neq 0\) 对一切 \(x \in U\) 成立。证明 \(Z_{U} \to U\) 是 U 的次数 \(\leq d\) 的覆盖（《一般拓扑学》第十一章），且 \(R \times U - Z_{U}\) 的连通分支个数 \(\leq d + 1\) 。
\item
  设 \(X'\) 是 X 中使 \(\Delta\) \(\neq 0\) 的点组成的集合。假设 \(X'\) 在 X 中稠密。用 \(\mathcal{A}(\text{resp. } \mathcal{B})\) 表示 \(X'\) （相应地 \(R \times X - Z_{X}\) ）的连通分支之集。证明
\end{enumerate}

\[
\operatorname{Card} (\mathscr {B}) \leq (d + 1) \operatorname{Card} (\mathscr {A}) \quad (\text { use   } b)).
\]

\begin{enumerate}
\def\labelenumi{\alph{enumi})}
\setcounter{enumi}{3}
\tightlist
\item
  假设 X 连通，且 \(d \geq 1\) 。证明
\end{enumerate}

\[
\operatorname{Card} (\mathcal {B}) \leq 1 + d \operatorname{Card} (\mathcal {A}).
\]

¶ 2) 设 V 是有限 n 维实向量空间，F 是 V 上的 d 次多项式函数。设 \(V'\) 是 V 中使 \(F \neq 0\) 的点组成的集合。证明 \(V'\) 的连通分支个数有限，且被一个只依赖于 n 与 d 的常数所界。（对 n 作归纳论证。化归到 F 无重因子的情形，并证明 V 可分解为 \(R \times X\) ，使习题 1 的结果可应用于 F。）\(^{10}\)

